% Generated by scripts/assemble.py; corrections belong in source/corrections.json.
\examhead{110}{安徽大学 2024 年数学分析试卷}{108}
% p109-b01; source PDF page 109
\par\Needspace{5\baselineskip}
\blockanchor{p109-b01}
\majorquestion{一}(15 分) 计算极限.\par
% p109-b02; source PDF page 109
\blockanchor{p109-b02}
(1) (7 分)  $\lim_{n\to\infty}n\left[\frac{1^{2}+3^{2}+\cdots+(2n+1)^{2}}{n^{3}}-\frac{4}{3}\right]$ .\par
% p109-b03; source PDF page 109
\blockanchor{p109-b03}
(2) (8 分) $\lim_{x \to \infty} \left[ \frac{1}{e} \left( 1 + \frac{1}{x} \right)^x \right]^x$.\par
% p109-b04; source PDF page 109
\par\Needspace{5\baselineskip}
\blockanchor{p109-b04}
\majorquestion{二}(15 分) 设数列 $\{x_{n}\}$ 有界, 且 $\lim_{n\to \infty}(x_{n + 1} - x_n) = 0$. 再设 $m = \varliminf_{n\to \infty}x_n$ 和 $M = \varlimsup_{n\to \infty}x_n$, 其中 $m < M$. 证明区间 $(m,M)$ 上任意一个数都是此数列的一个子列的极限.\par
% p109-b05; source PDF page 109
\par\Needspace{5\baselineskip}
\blockanchor{p109-b05}
\majorquestion{三}(20 分) 设函数 $f(x)$ 在区间 $I$ 上有定义, 令 $\omega_f(\delta) = \sup_{\substack{x,y \in I \\ |x - y| < \delta}} |f(x) - f(y)|$. 证明:\par
% p109-b06; source PDF page 109
\blockanchor{p109-b06}
(1) $\lim_{\delta \to 0^{+}}\omega_{f}(\delta)$ 存在.\par
% p109-b07; source PDF page 109
\blockanchor{p109-b07}
(2) $f(x)$ 在区间 $I$ 上一致连续等价于 $\lim_{\delta \to 0^{+}}\omega_{f}(\delta) = 0.$\par
% p109-b08; source PDF page 109
\par\Needspace{5\baselineskip}
\par\Needspace{6\baselineskip}
\blockanchor{p109-b08}
\majorquestion{四}(20 分) 设函数 $f(x)$ 为 $(a, b)$ 上的凸函数, 即对任意的 $x_1, x_2 \in (a, b)$ 以及 $\lambda \in (0, 1)$, 有\par
\nopagebreak[4]
% p109-b09; source PDF page 109
\blockanchor{p109-b09}
\displayformula{f (\lambda x _ {1} + (1 - \lambda) x _ {2}) \leq \lambda f (x _ {1}) + (1 - \lambda) f (x _ {2}).}
% p109-b10; source PDF page 109
\blockanchor{p109-b10}
证明：\par
% p109-b11; source PDF page 109
\blockanchor{p109-b11}
(1) 对任意的 $x \in (a, b)$, 左右导数 $f_{-}^{\prime}(x), f_{+}^{\prime}(x)$ 均存在, 且 $f_{-}^{\prime}(x) \leq f_{+}^{\prime}(x)$.\par
% p109-b12; source PDF page 109
\blockanchor{p109-b12}
(2) $f_{-}^{\prime}(x), f_{+}^{\prime}(x)$ 均在 $(a, b)$ 上单调递增.\par
% p109-b13; source PDF page 109
\par\Needspace{5\baselineskip}
\blockanchor{p109-b13}
\majorquestion{五}(20 分) 设 $f(n) = a_0 + \sum_{k=1}^{\infty} \frac{a_k}{n^k}$, 且满足 $|a_k| \leq M$, 这里 $n, k$ 均为正整数. 试证: 数项级数 $\sum_{n=1}^{\infty} f(n)$ 收敛的充要条件为 $a_0 = a_1 = 0$.\par
% p109-b14; source PDF page 109
\par\Needspace{5\baselineskip}
\blockanchor{p109-b14}
\majorquestion{六}(20 分) 给定方程 $x^{2} + y + \sin(xy) = 0$.\par
% p109-b15; source PDF page 109
\blockanchor{p109-b15}
(1) 说明在点 $(0,0)$ 的充分小的邻域内, 此方程确定唯一的可导的函数 $y = y(x)$, 使得 $y(0) = 0$, 并求出 $y = y(x)$ 的导函数表达式.\par
% p109-b16; source PDF page 109
\blockanchor{p109-b16}
(2) 判断在点 $(0,0)$ 的充分小的邻域内, 此方程是否确定唯一的函数 $x = x(y)$, 使得 $x(0) = 0$, 说明理由.\par
% p109-b17; source PDF page 109
\par\Needspace{5\baselineskip}
\blockanchor{p109-b17}
\majorquestion{七}(20 分) 计算 Gauss 曲面积分 $I = \iint_{S} \frac{\cos(\widehat{n,r})}{r^2} \, \mathrm{d}S$, 其中 $S$ 为光滑封闭曲面, 原点不在 $S$ 上, $r$ 为 $S$ 上动点至原定的距离, $(\widehat{n,r})$ 为动点处外法向量 $n$ 与径向 $r$ 的夹角.\par
% p109-b18; source PDF page 109
\par\Needspace{5\baselineskip}
\blockanchor{p109-b18}
\majorquestion{八}(20 分) 计算 $I(x) = \int_{0}^{+\infty} e^{-t^2} \cos(2xt) \, \mathrm{d}t$.\par
